\documentclass{beamer}

\newcommand{\week}{Banks 6 of 6}

\title{Bank Failures}
\subtitle{Mishkin Chapter 10}
\author{Econ 357}
\date{\week}

% Reference the shared preamble
\input{preamble}

% Set the footer -- change 
\setbeamertemplate{footline}{
  \leavevmode%
  \vspace{2ex}
  \hbox{%
    % Left box: Econ 457
    \begin{beamercolorbox}[wd=.4\paperwidth,ht=2.5ex,dp=1ex,left]{author in head/foot}%
      \hspace{1em}Econ 357
    \end{beamercolorbox}%
    % Middle box: Week
    \begin{beamercolorbox}[wd=.2\paperwidth,ht=2.5ex,dp=1ex,center]{date in head/foot}%
      \centering\week
    \end{beamercolorbox}%
    % Right box: Slide numbers
    \begin{beamercolorbox}[wd=.4\paperwidth,ht=2.5ex,dp=1ex,center]{date in head/foot}%
      \hfill\insertframenumber{} 
    \end{beamercolorbox}%
  }%
  \vskip0pt%
}

\begin{document}

\frame{\titlepage}

\begin{frame}
    \frametitle{Outline for Banks}
        \begin{enumerate}
            \item Role of Banks
            \item Bank Balance Sheets
            \item Bank Earnings
            \item Bank Management (and Investment Banks)
            \item Bank Failures
            \item \fbox{Bank Regulation}
        \end{enumerate}
\end{frame}


\begin{frame}
    \frametitle{Outline for Today's Lecture}
    
    \begin{enumerate}
        \item Moral Hazard and Too Big to Fail
        \item Regulation after the Great Depression
        \begin{itemize}
            \item Glass-Steagall (1933)
            \item Securities and Exchange Commission (1934)
        \end{itemize}
        \item Regulation after the Great Recession
        \begin{itemize}
            \item Stress Tests (2009)
            \item Dodd-Frank (2010)
        \end{itemize}
    \end{enumerate}

\end{frame}


\begin{frame}
\frametitle{Notes}

Notes:

    \begin{enumerate}
        \item In the interest of time, we are going to skip the discussion in Mishkin about bank consolidation and international banks.   You can skip Mishkin sections 11.3, 11.4, and 11.7.
        \item Mishkin covers the Great Financial Crisis (GFC) of 2008 in Chapter 12.   We will cover the GFC in depth later in the course.   At this point, we are just discussing the regulation that was enacted after the GFC.  Please read Mishkin sections 12.4 and 12.5.
    \end{enumerate}

\end{frame}

\begin{frame}[t]
\frametitle{1. Moral Hazard and Too Big to Fail}
\framesubtitle{}

    Moral hazard is a standard problem in insurance markets.   Insurers may be concerned that providing insurance may incentivize excessive risk taking.
    \vspace{1em}
    
    What do insurers do about this?
    
    \begin{itemize}
        \item Charge different prices based on riskiness of insured (their ability to do this is limited by law, and also because of the Lemons problem and markets could break down).
        \item Monitor behavior, regularly reassess risks
        \item Generally, insurers don’t ever make payments above what was initially contracted.
    \end{itemize}

\end{frame}

\begin{frame}[t]
\frametitle{1. Moral Hazard and Too Big to Fail}
\framesubtitle{}

    \textbf{“Too Big to Fail”}
    \vspace{1em}
    
    \begin{itemize}
        \item During a crisis, it is in the public’s best interest for the government to stop the bank run.   That may require providing insurance ex post, or even bailing out banks, for example by buying assets above “fire sale” prices.
        \item The government may do this, even if there was no insurance contract in place prior to the crisis.
    \end{itemize}
    \vspace{1em}
    
    The 2008 crisis saw many examples of an expansion of insurance, or what were commonly called  “bailouts”

\end{frame}

\begin{frame}[t]
\frametitle{1. Moral Hazard and Too Big to Fail}
\framesubtitle{}

    Moral Hazard / Too Big to Fail concerns may be especially acute in banking:
    
    \begin{itemize}
        \item Depositors have no incentive to monitor risks, because they are covered by FDIC guarantee. Riskier banks don’t have to pay more for deposit funding.
        \item Returns for bank owners/managers are very asymmetric.  They benefit significantly from the upside, where as the downside may be borne by taxpayers.
    \end{itemize}

\end{frame}

\begin{frame}[t]
\frametitle{1. Moral Hazard and Too Big to Fail}
\framesubtitle{}

    \textbf{FDIC Insurance}
    
    \begin{itemize}
        \item All banks are charged fees for being covered by FDIC insurance.
        \item Bank fees are adjusted based on perceived riskiness of the bank.
        \item The fees cover the costs to the taxpayers when the FDIC takes over a bank.   As a result, FDIC generally doesn’t cost the taxpayers money.
        \item However, these fees have not fully addressed the concerns about Moral Hazard.
    \end{itemize}

\end{frame}

\begin{frame}[t]
\frametitle{1. Moral Hazard and Too Big to Fail}
\framesubtitle{}

    \textbf{2008 Bailouts}
    \begin{itemize}
        \item Many banks and insurers received support from the federal government in 2008.
        \item The costs of the bailout in 2008 ended up being very modest, even zero in some calculations.
        \item However, the perception was very damaging.   A clear example was that even after AIG was taken over by the government, many AIG traders were paid multimillion dollar salaries.
        \item Many believed that Wall Street received a bailout, while Main Street did not
    \end{itemize}

\end{frame}

\begin{frame}[t]
\frametitle{1. Moral Hazard and Too Big to Fail}
\framesubtitle{}

    There was a lot of public outrage at banks and at the government after 2008.  
    \vspace{1em}
    
    Question: Was the government justified in bailing out banks?
    \vspace{1em}
    
    Discussion of Matt Taibbi and Robert Shiller
    \vspace{1em}

\end{frame}

\begin{frame}[t]
\frametitle{2. Regulation after the Great Depression}
\framesubtitle{}

\textbf{Glass-Steagall Act of (1933)}

\begin{itemize}
    \item Created the FDIC
    \item Separation of Commercial Banks and Investment Banks
    \item Commercial banks: take deposits and make loans
    \item Investment banks: underwrote and deal in securities
    \item Intended to make the financial system safer by splitting the riskier investment banking activities from the deposit-based commercial banking activities.
    \item For example, The House of Morgan was split into Morgan Stanley (Investment Bank) and JP Morgan (Commercial Bank).
\end{itemize}

\end{frame}

\begin{frame}[t]
\frametitle{2. Regulation after the Great Depression}
\framesubtitle{}

    \textbf{Post-script}
    The FDIC remains a very important part of our financial regulatory system.  In fact, it has somewhat expanded it's ability to guarantee short term creditors in banks and it played a very large role in providing guarantees during the 2008 financial crisis.
    \vspace{1em}
    
    In contrast, Glass-Steagall was repealed in 1999 by Gramm-Leach-Bliley.   After that, commercial banks and investment banks were allowed to merge.  This contributed to a trend of consolidation in the banking industry, as can be seen in the next few slides.
    \vspace{1em}

\end{frame}

\begin{frame}[t]
\frametitle{2. Regulation after the Great Depression}
\framesubtitle{Bank Consolidation}

    \centering
    \textbf{Number of Commercial Banks in the US, 1934-2020}
    \vspace{1em}
    
    \includegraphics[width=0.8\textwidth]{Econ357/figures/03_banks_6_consolidation.png}
    
\end{frame}


\begin{frame}[t]
\frametitle{2. Regulation after the Great Depression}
\framesubtitle{Bank Concentration in 2025}

\footnotesize
\begin{table}
\centering
\begin{tblr}{ 
    colspec = { Q[l,wd=5cm] Q[c,wd=2cm] Q[c,wd=2cm]} 
}
\toprule
Bank &  Assets (\$ tr) &  Share of All Assets (\%) \\
\midrule
1. J.P. Morgan Chase \& Co. & 2.74 & 15.4 \\
2. Bank of America Corp. & 2.38 & 13.4 \\
3. Citigroup Inc. & 1.96 & 11.0 \\
4. Wells Fargo \& Co. & 1.89 & 10.6 \\
5. Goldman Sachs Group &  0.93 &  5.2 \\
6. Morgan Stanley & 0.88 &  4.9 \\
7. U.S. Bankcorp & 0.48  & 2.7 \\
8. P N C Financial Services Group & 0.39 & 2.2 \\
9. T D Group U S & 0.38 & 2.1 \\
10. Capital One Financial Corp. & 0.37 & 2.1\\
\midrule
Total & 12.40 & 69.6\\
\bottomrule
\end{tblr}
\end{table}
\end{frame}

\begin{frame}[t]
\frametitle{2. Regulation after the Great Depression}
\framesubtitle{}

    \textbf{Discussion question}: Did the end of Glass-Steagall lead to the financial crisis in 2008?   It's a complicated question with no clear answer.
    \vspace{1em}
    
    On the one hand, the investment side of banks took a lot of risk to generate fees.   At some point this endangered depositors, and the government was forced to step in.   That was precisely what Glass-Steagall had sought to prevent.
    \vspace{1em}
    
    On the other hand, the risk taking likely would have happened anyways, and the investment banks may have required bailouts (the federal government took over AIG, which was an insurer).   Moreover, its not clear that the housing bubble, which was the root cause of the crisis, had anything to do with Glass-Steagall.


\end{frame}


\end{document}